% Macros and settings used by both the pdf and the web version.

% Theorem-like environments, numbered within chapters. The nodes of the dependency graph are those
% listed in the `thms` option of web.tex; remarks, the question and the conjecture are not nodes.
\newtheorem{theorem}{Theorem}[chapter]
\newtheorem{proposition}[theorem]{Proposition}
\newtheorem{lemma}[theorem]{Lemma}
\newtheorem{corollary}[theorem]{Corollary}
\newtheorem{facts}[theorem]{Facts}
\newtheorem{conjecture}[theorem]{Conjecture}
% Statements of the 1992 paper that Part I formalizes in its own form (Appendix A). Part II
% states them with their numbers and points to their nodes in Part I; they are not nodes.
\newtheorem{sharedfacts}[theorem]{Facts}
\newtheorem{sharedlemma}[theorem]{Lemma}

\theoremstyle{definition}
\newtheorem{definition}[theorem]{Definition}
\newtheorem{remark}[theorem]{Remark}
\newtheorem{question}[theorem]{Question}

% Chapters 1--7 are the sections of the paper, and their statements carry the paper's numbers.
% The supporting statements of a chapter, which the paper does not state, come first, in a
% section opened by \supporting, and are numbered with letters (Lemma 5.A, ...); \paperitems
% then restarts the arabic numbering.
\newcommand{\supporting}{\renewcommand{\thetheorem}{\thechapter.\Alph{theorem}}}
\newcommand{\paperitems}{%
  \setcounter{theorem}{0}\renewcommand{\thetheorem}{\thechapter.\arabic{theorem}}}

% Parts are numbered I and II (plasTeX numbers them in arabic by default).
\renewcommand{\thepart}{\Roman{part}}

% Part II follows the papers of 1992 and of Li. Its chapters are numbered II.1, II.2, ..., so
% that the statements of the 1992 paper keep their numbers behind the prefix: Lemma II.1.4 is
% Lemma 1.4 of that paper. \secondpart also undoes the effect of \appendix on the chapter names.
\newcommand{\secondpart}{%
  \setcounter{chapter}{0}\renewcommand{\thechapter}{II.\arabic{chapter}}\restorechapters}

% Notation of the paper.
\newcommand{\QQ}{\mathbb{Q}}
\newcommand{\ZZ}{\mathbb{Z}}
\newcommand{\NN}{\mathbb{N}}
\newcommand{\FF}{\mathbb{F}}
\newcommand{\Qbar}{\overline{\QQ}}
\newcommand{\Om}{\Omega}
\newcommand{\Ct}[1]{[C_2]^{#1}}       % iterated wreath power of C_2
\newcommand{\Gal}{\operatorname{Gal}}
\newcommand{\Hom}{\operatorname{Hom}}
\newcommand{\ord}{\operatorname{ord}}
\newcommand{\sgn}{\operatorname{sign}}
\newcommand{\dvd}{\mathrel{\mid}}
\newcommand{\rad}{\operatorname{rad}}
\newcommand{\val}{\operatorname{val}}
\newcommand{\Leg}[2]{\left(\frac{#1}{#2}\right)}
\newcommand{\sq}{\square}
\newcommand{\cA}{\mathcal{A}}
\newcommand{\eps}{\varepsilon}

% Marks a place where the formal statement deviates from the paper.
\newcommand{\formalnote}[1]{\par\smallskip\noindent\textit{Formalization:} #1\par\smallskip}
